\src{249}
\seminar{Năm thứ 13, 1960/61, số 221\\Tháng 5 năm 1961}
\paperhead{CÁC KỸ THUẬT XÂY DỰNG VÀ CÁC ĐỊNH LÝ TỒN TẠI TRONG HÌNH HỌC ĐẠI SỐ\\
IV. CÁC LƯỢC ĐỒ HILBERT}%
  {Alexander GROTHENDIECK}

\subsection*{Mở đầu}
\addcontentsline{toc}{subsection}{Mở đầu}

Các kỹ thuật được trình bày trong [2], I và II, về cơ bản không phụ thuộc vào
bất kỳ giả thiết xạ ảnh nào đối với các lược đồ được xét. Tiếc rằng hiện nay
chúng chưa cho phép giải các bài toán tồn tại được đặt ra trong Exposé II của
[2]. Trong exposé này và exposé kế tiếp, chúng tôi sẽ giải các bài toán ấy dưới
các giả thiết xạ ảnh. Các kỹ thuật được sử dụng mang tính xạ ảnh điển hình và
hầu như không dùng đến các kết quả của I và II trong [2]. Ở đây, chúng tôi sẽ
xây dựng các «lược đồ Hilbert», nhằm thay thế việc sử dụng các tọa độ Chow, như
đã nói trong [2], II, số 2. Trong exposé kế tiếp, lý thuyết phép lấy thương đối
với các lược đồ được phát triển trong III, kết hợp với lý thuyết các lược đồ
Hilbert, sẽ cho phép chúng tôi, chẳng hạn, xây dựng các lược đồ Picard (được
định nghĩa trong [2], II, số 3) dưới những điều kiện khá tổng quát.

Tóm lại, có thể nói rằng hiện nay ta đã có một kỹ thuật xây dựng xạ ảnh khá
thỏa đáng, ngoại trừ việc vẫn còn thiếu\footnote{Xem phần bổ sung ở cuối
exposé.} một định lý về phép lấy thương theo các nhóm như nhóm xạ ảnh, tác động
«không có điểm bất động» (xem [2], III, số 8). Tình hình dường như còn thuận lợi
hơn đôi chút trong hình học giải tích (khi ta chỉ xét các không gian giải tích
«xạ ảnh» trên một không gian giải tích đã cho), bởi vì đối với các không gian
giải tích, khó khăn của phép lấy thương theo một nhóm tác động tốt sẽ biến mất.
Mặt khác, trong hình học đại số cũng như trong hình học giải tích, vẫn còn phải
phát triển một kỹ thuật xây dựng có hiệu lực mà không cần các giả thiết xạ ảnh.

\fsection{1}{Các họ bị chặn của các bó: các tính chất ổn định}

Cho $k$ là một trường và $X$ là một $k$-tiền lược đồ, giả sử thuộc kiểu hữu hạn
để đơn giản. Với mỗi mở rộng $K/k$, ta do đó thu được một $K$-tiền lược đồ
$X_K=X\otimes_k K$. Nếu $\mathcal F$ là một bó coherent trên $X_K$ và nếu
$K'$ là một mở rộng của $K$, thì
$\mathcal F\otimes_K K'=\mathcal F_{K'}$ là một bó quasi-coherent trên
$X_K\otimes_K K'=X_{K'}$. Như vậy, nếu $K$ và $K'$ là hai mở rộng tùy ý của
$k$, $\mathcal F$ là một bó quasi-coherent trên $X_K$ và $\mathcal F'$
\src{250}
là một bó quasi-coherent trên $X_{K'}$, ta nói $\mathcal F$ và
$\mathcal F'$ là \emph{tương đương} nếu tồn tại các $k$-đồng cấu từ $K$,
$K'$ vào cùng một mở rộng $K''$ của $k$ sao cho $\mathcal F_{K''}$ và
$\mathcal F'_{K''}$ đẳng cấu trên $X_{K''}$. Đây là một quan hệ tương đương;
ta sẽ quan tâm đến các lớp tương đương của các bó theo quan hệ này và đến các
tập hợp lớp tương đương. Lưu ý rằng, nếu $X_0$ thuộc kiểu hữu hạn trên
$k$, thì mỗi lớp bó coherent có thể được xác định bởi một bó coherent trên
một $X_K$, trong đó $K$ là một mở rộng hữu hạn sinh của $k$. Vì vậy, trong
định nghĩa các lớp bó coherent, ta có thể chỉ xét các mở rộng đóng đại số của
$k$; cũng có thể chỉ xét một mở rộng đóng đại số cố định $\Omega$ của $k$, có
bậc siêu việt vô hạn. Khi ấy, hai bó coherent $\mathcal F$ và $\mathcal F'$
trên $X_\Omega$ tương đương khi và chỉ khi tồn tại một $k$-tự đẳng cấu $\sigma$
của $\Omega$ sao cho $\mathcal F\otimes_\Omega(\Omega,\sigma)$ đẳng cấu với
$\mathcal F'$. Do đó có một tương ứng song ánh giữa các lớp bó coherent theo
một trong hai định nghĩa.

Cho $E$, $E'$ là hai tập hợp các lớp bó coherent trên $X$. Xét các lớp của mọi
bó có dạng $\mathcal F\otimes\mathcal F'$, trong đó $\mathcal F$ và
$\mathcal F'$ là các bó coherent trên cùng một $X_K$, lớp của $\mathcal F$
thuộc $E$ và lớp của $\mathcal F'$ thuộc $E'$. Ta thu được một tập hợp các lớp
bó coherent, ký hiệu là $E\otimes E'$. Tương tự, ta định nghĩa các
$\operatorname{Tor}_i(E,E')$, v.v. Tổng quát hơn, với mỗi hàm $U$ gán cho mọi
dãy $\mathcal F_1,\ldots,\mathcal F_n$ gồm $n$ bó coherent trên cùng một $X_K$
một tập hợp $U(\mathcal F_1,\ldots,\mathcal F_n)$ các bó coherent trên $X_K$,
và có tính chất tương thích hiển nhiên với các đẳng cấu bó và các ảnh ngược qua
phép đổi cơ sở, ta liên kết một hàm, cũng ký hiệu là $U$, gán cho mọi dãy gồm
$n$ tập hợp $E_1,\ldots,E_n$ các lớp bó coherent một tập hợp
$U(E_1,\ldots,E_n)$ các lớp bó coherent.

Mục đích của chúng ta trong mục này là định nghĩa một số tập hợp lớp bó, gọi là
\emph{bị chặn}, và chứng minh rằng các phép toán $U$ thông dụng nhất, khi áp
dụng cho các tập hợp bị chặn, vẫn cho các tập hợp bị chặn.

Cho $X$ là một tiền lược đồ thuộc kiểu hữu hạn trên một lược đồ Noether $S$.
Với mỗi $s\in S$, thớ $X_s$ là một tiền lược đồ thuộc kiểu hữu hạn trên
$k(s)$, và ta sẽ xét các lớp bó coherent trên $X_s$ theo nghĩa trên. Điều này
mang lại nghĩa cho cụm từ: lớp bó coherent trên một thớ của $X/S$, cũng như
cho các cụm từ tương tự. Cũng vậy, khi tiến hành riêng trên từng thớ, ta vẫn có
thể xét các phép toán như $E\otimes E'$, v.v., gán cho mỗi hệ gồm những tập hợp
\src{251}
các lớp bó coherent trên các thớ của $X/S$ một tập hợp khác gồm các lớp bó
coherent trên các thớ của $X/S$.

\result{ĐỊNH NGHĨA}{1.1} Cho $E$ là một tập hợp các lớp bó coherent trên các
thớ của $X/S$. Ta nói $E$ là \emph{bị chặn} nếu tồn tại một tiền lược đồ $S'$
thuộc kiểu hữu hạn trên $S$ và một bó coherent $\mathcal F'$ trên
$X'=X\times_S S'$ sao cho $E$ được chứa trong tập hợp các lớp bó trên các thớ
của $X/S$ do $\mathcal F'$ xác định.

Theo định nghĩa, tập hợp cuối cùng này gán cho mỗi $s\in S$ các lớp của những bó
$\mathcal F'\otimes_{S'}k(s')$, trong đó $s'$ chạy qua các điểm của $S'$ nằm
trên $s$ (do đó $k(s')$ là một mở rộng của $k(s)$, và
$X_s\otimes_{k(s)}k(s')$ được đồng nhất với thớ
$X'\otimes_{S'}k(s')=X'_{s'}$ của $X'$ tại $s'$). Có thể nói rằng các họ bị
chặn là những họ được chứa trong một \emph{họ đại số} các bó coherent, được
tham số hóa bởi một $S'$ thuộc kiểu hữu hạn trên $S$.

Một hợp hữu hạn của các họ bị chặn là bị chặn (lấy tiền lược đồ tổng rời nhau
của các tiền lược đồ tham số $S_i$ xác định những họ đại số bao chứa). Một phép
đổi cơ sở $T\longrightarrow S$ biến một họ bị chặn tương đối với $X/S$ thành
một họ bị chặn tương đối với $X_T/T$; mệnh đề đảo đúng nếu
$T\longrightarrow S$ là toàn ánh (hoặc, tổng quát hơn, nếu ảnh của nó chứa các
$s$ thực sự xuất hiện trong họ $E$ đã cho đối với $X/S$). Về lý thuyết, điều
này quy việc xác định các họ bị chặn về trường hợp $S$ là phổ của một đại số
thuộc kiểu hữu hạn trên vành số nguyên $\mathbf Z$.

Nếu $E$ và $E'$ là các họ bị chặn gồm các lớp bó tương đối với $X/S$, thì
$E\otimes E'$ cũng bị chặn. Thật vậy, giả sử $E$ và $E'$ được chứa tương ứng
trong các họ đại số xác định bởi $T\longrightarrow S$ và $\mathcal F$ trên
$X_T$, bởi $T'\longrightarrow S$ và $\mathcal F'$ trên $X_{T'}$. Khi đó,
$E\otimes E'$ được chứa trong họ đại số xác định bởi
$T''=T\times_S T'\longrightarrow S$ và bó $\mathcal F''$ trên $X_{T''}$, là
tích tensor của các ảnh ngược của $\mathcal F$ và $\mathcal F'$ trên
$X_{T''}$. Lập luận này chỉ đúng bởi vì hàm tử
$\mathcal F\otimes\mathcal F'$ theo các biến $\mathcal F$, $\mathcal F'$ là
khớp phải, nên giao hoán với mở rộng cơ sở (và đặc biệt với phép chuyển sang
các thớ). Nó không áp dụng trực tiếp cho các phép toán địa phương như
$\operatorname{Tor}_i(E,E')$, $\operatorname{Hom}(E,E')$,
$\operatorname{Ext}^i(E,E')$. Tuy nhiên, có thể chứng minh rằng các phép toán
này vẫn biến các tập hợp bị chặn thành các tập hợp bị chặn, bằng cách tiến hành
như đối với $E\otimes E'$, nhưng sử dụng thêm các kết quả thuộc kiểu sau đây
(tất cả đều có trong [3], IV, 6.11): một họ bị chặn $E$ luôn được chứa trong
một họ đại số xác định bởi một bó coherent $\mathcal F$ trên một $X_T$ ($T$
thuộc kiểu hữu hạn trên $S$), phẳng trên $T$. (Ta «cắt thành từng mảnh» không
gian tham số ban đầu.) Những
\src{252}
tính chất phẳng của các bó thích hợp quả thực bảo đảm rằng các phép toán như
$\operatorname{Tor}_i(\mathcal F,\mathcal F')$ giao hoán với mọi phép đổi cơ
sở. Cùng phương pháp ấy áp dụng cho các phép toán mang tính toàn cục: ảnh trực
tiếp và ảnh trực tiếp dẫn xuất của các bó coherent qua các cấu xạ riêng, các
$\operatorname{Ext}$ toàn cục tương đối với các cấu xạ riêng (xem [5], III,
\S~6), v.v.; mọi phép toán này đều biến các họ bị chặn của những bó thành các
họ bị chặn của những bó. (Lưu ý. -- Ở đây những tiền lược đồ mà trên đó ta lấy
các bó khác nhau có thể thay đổi qua các phép toán thuộc kiểu đang xét.)

Hai phát biểu sau được chứng minh về cơ bản bằng cùng kỹ thuật về tính phẳng;
đối với sự phân tích nguyên sơ trên các thớ của một cấu xạ thuộc kiểu hữu hạn,
xem riêng [5], IV.

\result{MỆNH ĐỀ}{1.2} Cho $E$, $E'$ là những tập hợp bị chặn gồm các lớp bó
trên các thớ của $X/S$, giả sử $X$ riêng trên $S$. Khi đó:

\medskip
\noindent (i) Họ các hạt nhân, đối hạt nhân và ảnh của những đồng cấu
$\mathcal F\longrightarrow\mathcal F'$, trong đó lớp của $\mathcal F$ thuộc
$E$ và lớp của $\mathcal F'$ thuộc $E'$, là bị chặn.

\smallskip
\noindent (ii) Họ các bó $\mathcal F''$ là những mở rộng của một $\mathcal F$
bởi một $\mathcal F'$, trong đó lớp của $\mathcal F$ thuộc $E$ và lớp của
$\mathcal F'$ thuộc $E'$, là bị chặn.

Sau một phép đổi cơ sở thích hợp, có thể giả sử $E$ và $E'$ lần lượt được xác
định bởi một bó coherent $\mathcal G$ và một bó coherent $\mathcal G'$ trên
một $X_T/T$, với $T$ thuộc kiểu hữu hạn trên $S$. Hơn nữa, có thể giả sử một
số giả thiết về tính phẳng được thỏa mãn, kéo theo việc tạo thành các bó
$f_{T*}(\operatorname{Hom}_{X_T}(\mathcal G,\mathcal G'))$ và
$\operatorname{Ext}^{1}_{f_T}(\mathcal G,\mathcal G')$ giao hoán với phép đổi
cơ sở theo một cấu xạ tùy ý $T'\longrightarrow T$. Cũng có thể giả sử các bó
coherent nói trên là tự do địa phương trên $T$. Khi đó, gọi $T_0$ và $T_1$
là các phân thớ vectơ trên $T$ mà bó các mầm tiết diện của chúng lần lượt là
các bó nói trên. Khi ấy ta xác định chính tắc một đồng cấu
$\mathcal G_{T_0}\longrightarrow\mathcal G'_{T_0}$ giữa các bó coherent trên
$X_{T_0}$, và một mở rộng
\[
0\longrightarrow\mathcal G_{T_1}\longrightarrow\mathcal G''
\longrightarrow\mathcal G'_{T_1}\longrightarrow0
\]
gồm các bó coherent trên $X_{T_1}$, có một tính chất phổ dụng hiển nhiên. Bó
thứ hai này xác định một họ đại số bao chứa họ đang xét trong (ii). Điều tương
tự cũng đúng với hạt nhân, đối hạt nhân và ảnh của đồng cấu
\src{253}
trước đó, mà việc xét nó chứng minh (i) (miễn là giả sử đối hạt nhân phẳng trên
$T_0$, trường hợp mà ta lại quy về bằng cách cắt $T_0$ thành từng
mảnh\ldots).

\result{MỆNH ĐỀ}{1.3} Cho $E$ là một họ bị chặn gồm các lớp bó trên các thớ
của $X/S$. Khi đó, các lớp của các bó cấu trúc của
$(\operatorname{supp}\mathcal F)_{\mathrm{red}}$, trong đó $\mathcal F$ là một
bó coherent trên một $X_K$, với $K$ đóng đại số, và lớp của nó thuộc $E$, tạo
thành một họ bị chặn.

Ở đây, $(\operatorname{supp}\mathcal F)_{\mathrm{red}}$ ký hiệu cho giá của
$\mathcal F$ được trang bị cấu trúc rút gọn cảm sinh; tức là bó cấu trúc của nó
là thương của $\mathcal O_{X_K}$ theo bó iđêan lớn nhất xác định
$\operatorname{supp}\mathcal F$. Có thể chứng minh kết quả tương tự (1.3) cho
những bó được suy ra chính tắc từ $\mathcal F$ bằng lý thuyết phân tích nguyên
sơ, chẳng hạn các $\mathcal F/\mathcal F_i$, trong đó các $\mathcal F_i$ là
những bó con của $\mathcal F$ nguyên sơ ứng với các thành phần của giá của
$\mathcal F$ và tối tiểu đối với tính chất này; hoặc các
$\mathcal O_{X_K}/\mathfrak p$, trong đó $\mathfrak p$ là một bó iđêan nguyên
tố liên kết với $\mathcal F$; hoặc các $\mathcal O_{X_K}/\mathfrak q$, trong đó
$\mathfrak q$ là một bó iđêan nguyên sơ liên kết với một thành phần của giá của
$\mathcal F$ (trường tham chiếu được giả sử đóng đại số).

\fsection{2}{Các họ bị chặn và đa thức Hilbert}

Trong phần tiếp theo, ta giả sử $X$ xạ ảnh trên $S$ và được trang bị một bó rất
rộng, ký hiệu $\mathcal O_X(1)$. Với mọi mở rộng $K$ của trường thặng dư $k(s)$
tại một điểm $s$ của $S$, ta xét trên $X_K$ bó tương ứng
$\mathcal O_{X_K}(1)$, bó này vẫn rất rộng.

Với mỗi bó coherent $\mathcal F$ trên $X_K$, ta gán hàm
\[
P_{\mathcal F}(n)=\text{đặc trưng Euler--Poincaré của $\mathcal F(n)$
trên $X_K$},
\]
là một đa thức theo số nguyên $n$, gọi là \emph{đa thức Hilbert} của
$\mathcal F$. Với các giá trị đủ lớn của $n$, $P(n)$ không gì khác hơn là số
chiều của $H^0(X_K,\mathcal F(n))$ trên $K$, vì các
$H^i(X_K,\mathcal F(n))$ bằng không khi $i>0$ và $n$ đủ lớn.

Bây giờ, nếu $\mathcal F$ là một bó coherent trên $X$, phẳng trên $S$, thì các
đa thức Hilbert của những bó $\mathcal F_s$ cảm sinh trên các thớ $X_s$ thuộc
cùng một thành phần liên thông của $S$ đều bằng nhau [5], III, \S~7. Do đó
(không cần giả thiết phẳng), tập hợp các đa thức Hilbert của những bó
$\mathcal F_s$, $s\in S$, là hữu hạn đối với mọi bó coherent $\mathcal F$ trên
$X$.
\src{254}
Mặt khác, nhắc lại rằng nếu $\mathcal F$ là một bó coherent trên $X$, thì nó
đẳng cấu với một bó thương của bó $\mathcal O_X(-n)^N$, với $n$ và $N$ đủ
lớn. Do đó, các bó $\mathcal F_s$ cảm sinh trên các thớ cũng là thương của bó
$\mathcal O(-n)$ trên thớ.

Từ hai nhận xét này, ta suy ra chiều «chỉ khi» của định lý sau:

\result{ĐỊNH LÝ}{2.1} Cho $X$ xạ ảnh trên $S$ Noether, với
$\mathcal O_X(1)$ rất rộng trên $X$ tương đối với $S$. Cho $E$ là một tập hợp
các lớp bó trên các thớ của $X/S$. Để $E$ bị chặn, điều kiện cần là nó thỏa mãn các điều kiện sau:

\medskip
\noindent (a) Tồn tại một bó coherent $\mathcal L$ trên $X$ (có thể giả sử có
dạng $\mathcal O_X(-n)^N$) sao cho $E$ được chứa trong họ các lớp của những bó
coherent là thương của các bó dạng $\mathcal L_K$.

\smallskip
\noindent (b) Các đa thức Hilbert $P_{\mathcal F}$ của những bó $\mathcal F$
có lớp thuộc $E$ đều thuộc cùng một tập hợp hữu hạn các đa thức.

Còn phải chứng minh chiều «nếu», đây sẽ là một trường hợp riêng của một kết quả
chính xác hơn. Với mọi môđun coherent $\mathcal F$ trên một tiền lược đồ thuộc
kiểu hữu hạn trên một trường $K$, và mọi số nguyên $r$, gọi $N_r$ là môđun con
của $\mathcal F$ mà các tiết diện trên một tập mở là những tiết diện của
$\mathcal F$ trên tập mở đó có giá mang chiều $<r$. Khi đó
$N_r=\mathcal F$ nếu $r>\dim\operatorname{supp}\mathcal F$, còn $N_r=0$ nếu
$r\leq0$, và ta thu được một lọc tăng hữu hạn của $\mathcal F$ mà các thương
$N_r/N_{r+1}$ có các chu trình nguyên tố
liên kết đúng bằng những chu trình nguyên tố liên kết với $\mathcal F$ mang
chiều $r$. Ta đặt
\[
\mathcal F_{(r)}=\mathcal F/N_r,
\]
thì các chu trình nguyên tố liên kết với $\mathcal F_{(r)}$ đúng bằng những chu
trình nguyên tố liên kết với $\mathcal F$ mang chiều $\geq r$; đặc biệt, nó
bằng $\mathcal F$ khi và chỉ khi các chu trình nguyên tố liên kết với
$\mathcal F$ mang chiều $\geq r$. Với ký hiệu này:

\result{ĐỊNH LÝ}{2.2} Dưới các giả thiết của (2.1), cho $s$ là một số nguyên
và giả sử $E$ thỏa mãn điều kiện (a) cùng dạng yếu hơn sau đây của (b):

\medskip
\noindent $(b_s)$ Các đa thức Poincaré $P_{\mathcal F}$ của những bó
$\mathcal F$ có lớp thuộc $E$ có các hệ số ở bậc $\leq s-1$ được chặn.

Dưới các điều kiện này, các bó $\mathcal F_{(s)}$ (khi lớp của $\mathcal F$
chạy trong $E$) tạo thành một họ bị chặn. Hơn nữa, các hệ số bậc $s-2$ của
các $P_{\mathcal F}$ bị chặn dưới.
\src{255}
Do đó:

\result{HỆ QUẢ}{2.3} Giả sử những bó $\mathcal F$ có lớp thuộc $E$ sao cho
mọi chu trình nguyên tố liên kết của chúng đều mang chiều $d$ thỏa mãn
$s\leq d\leq r$. Khi đó, trong (2.1) (b), chỉ cần xét các hệ số của
$P_{\mathcal F}$ ở các bậc từ $s-1$ đến $r$.

Phần còn lại của mục này dành cho việc phác thảo chứng minh (2.2). Các bổ đề
then chốt là hai bổ đề sau, trong đó bổ đề thứ nhất đã quen thuộc (và tóm tắt
nội dung toán học hữu dụng của các tọa độ Chow):

\result{BỔ ĐỀ 2.4 (CHOW)}{} Xét các bó cấu trúc của những lược đồ con $Y$ của
các thớ $X_K$ ($K$ chạy qua các mở rộng đóng đại số của những trường thặng dư
của $S$), với $Y$ rút gọn và mọi thành phần của nó cùng mang chiều $r$
($\mathcal O_Y$ được xem như một bó thương của $\mathcal O_X$). Nếu các bậc
của $Y$ bị chặn, thì các $Y$ tạo thành một họ bị chặn.

Ở đây, bậc $a$ của $Y$ có thể được định nghĩa thuận tiện nhất bằng hệ số của
số hạng đầu của
\[
P_{\mathcal O_Y}=a\frac{n^r}{r!}+\cdots.
\]

\result{BỔ ĐỀ}{2.5} Cho $\mathcal L$ là một bó coherent trên $X$ và $E$ là
một tập hợp các lớp bó thương $\mathcal F$ của những bó $\mathcal L_K$ ($K$
là một mở rộng trường thặng dư của $S$). Giả sử các thớ của $X$ trên $S$ mang
chiều $\leq r$, và đặt
\[
P_{\mathcal F}(n)=a_{\mathcal F}\frac{n^r}{r!}
+b_{\mathcal F}\frac{n^{r-1}}{(r-1)!}
+\text{các số hạng bậc $<r-1$}.
\]
Khi đó, hệ số $a_{\mathcal F}$ bị chặn, còn $b_{\mathcal F}$ bị chặn dưới.
Nếu $b_{\mathcal F}$ bị chặn, thì họ các $\mathcal F_{(r)}$ bị chặn.

\emph{Chứng minh.} -- Sau khi thay $S$, nếu cần, bằng một hợp các lược đồ con
của $S$ phủ $S$, ta có thể giả sử tồn tại một cấu xạ hữu hạn
$f:X\longrightarrow\mathbf P^r_S$ sao cho $\mathcal O_X(1)$ đẳng cấu với ảnh
ngược của $\mathcal O_{\mathbf P^r_S}(1)$; do đó, với mọi bó coherent
$\mathcal F$ trên $X$, ta có
$P_{\mathcal F}=P_{f_*(\mathcal F)}$. Mặt khác, dễ kiểm tra (bằng kỹ thuật của
mục trước) rằng một tập hợp các bó $\mathcal F$ trên $X$ bị chặn khi và chỉ khi
tập hợp các $f_*(\mathcal F)$ bị chặn. Cuối cùng, ta có
\[
f_*(\mathcal F)_{(r)}=f_*(\mathcal F_{(r)}).
\]
\src{256}
Như vậy, ta quy về trường hợp $X=\mathbf P^r_S$. Hơn nữa, có thể giả sử
$\mathcal L=\mathcal O_{\mathbf P^r_S}(k)^s$ với $k$, $s$ thích hợp. Hệ số
$a_{\mathcal F}$ thỏa mãn
\[
0\leq a_{\mathcal F}\leq s
\]
và do đó bị chặn. Mặt khác, việc các $P_{\mathcal F}(n)$ có hệ số của
$n^{r-1}$ bị chặn dưới (tương ứng, bị chặn) tương đương với việc khẳng định
điều đó cho
$P_{\mathcal F}(n-k)=P_{\mathcal F(-k)}(n)$. Như vậy, ta quy về trường hợp
\[
\mathcal L=\mathcal O_{\mathbf P^r_S}^{s}.
\]

Xét dãy khớp
\[
0\longrightarrow N_r\longrightarrow\mathcal F
\longrightarrow\mathcal F_{(r)}\longrightarrow0,
\]
từ đó
\[
P_{\mathcal F}=P_{\mathcal F_{(r)}}+P_{N_r},
\]
và vì hệ số của $n^{r-1}$ trong $P_{N_r}$ không âm (do
$\dim\operatorname{supp}N_r\leq r-1$), ta có
\[
b_{\mathcal F_{(r)}}\leq b_{\mathcal F}.
\]
Điều này cho phép ta, để chứng minh các khẳng định của bổ đề, thay
$\mathcal F$ bằng $\mathcal F_{(r)}$, tức là giả sử những bó thương
$\mathcal F$ của $\mathcal L$ đang xét là không xoắn.

Vì $\mathbf P^r_K$ chuẩn tắc, suy ra $\mathcal F$ tự do địa phương hạng
$a=a_{\mathcal F}$ trên một tập mở $U=\mathbf P^r_K-Y$, trong đó $Y$ có đối
chiều $\geq2$. Do đó, $\bigwedge^a\mathcal F$ là một bó trên
$\mathbf P^r_K$ mà hạn chế của nó lên $U$ là khả nghịch; vì vậy (do
$\mathbf P^r_K$ chính quy và $Y$ có đối chiều $\geq2$), hạn chế này đẳng cấu
với hạn chế của một bó khả nghịch trên $\mathbf P^r_K$, được xác định sai khác
một đẳng cấu. Bó sau có dạng $\mathcal O_{\mathbf P^r_K}(d)$ với một số nguyên
$d$ được xác định duy nhất. Vì $\bigwedge^a\mathcal F$ là thương của
$\bigwedge^a\mathcal O_{\mathbf P^r_K}^{n}\simeq
\mathcal O_{\mathbf P^r_K}^{N}$, với $N=\binom{s}{a}$, nó có $N$ tiết diện chính tắc, do đó xác
định các tiết diện của
\src{257}
$\mathcal O_{\mathbf P^r_K}(d)$ trên $U$, là những hạn chế của các tiết diện
$s_i$ ($1\leq i\leq N$) của $\mathcal O_{\mathbf P^r_K}(d)$ (vì
$\mathbf P^r_K$ chuẩn tắc và $Y$ có đối chiều $\geq2$). Các tiết diện này sinh
$\mathcal O_{\mathbf P^r_K}(d)$ tại các điểm của $U$, nên chúng không đồng thời
bằng không; điều này kéo theo $d\geq0$. Mặt khác, một phép tính đơn giản cho
\[
b_{\mathcal F}=a_{\mathcal F}(r+1)/2+d.
\]
Đặc biệt, điều này cho thấy $b_{\mathcal F}>0$, nên $b_{\mathcal F}$ bị chặn
dưới. Nó bị chặn khi và chỉ khi $d$ bị chặn; ta chứng minh rằng khi đó
$\mathcal F$ thuộc một họ bị chặn. Có thể giả sử $a_{\mathcal F}$ và
$b_{\mathcal F}$ cố định, lần lượt bằng $a$ và $b$, nên $d$ cố định. Việc biết
$N$ tiết diện $s_i$ của $\mathcal O_{\mathbf P^r_K}(d)$, tức là biết một đồng
cấu
\[
s:\bigwedge^a\mathcal L_K\longrightarrow\mathcal O_{\mathbf P^r_K}(d),
\]
cho phép khôi phục $\mathcal F$ như đối ảnh của đồng cấu hợp thành tương ứng:
\[
\mathcal L_K\longrightarrow
\operatorname{Hom}\!\left(\bigwedge^{a-1}\mathcal L_K,
\bigwedge^a\mathcal L_K\right)\longrightarrow
\operatorname{Hom}\!\left(\bigwedge^{a-1}\mathcal L_K,
\mathcal O_{\mathbf P^r_K}(d)\right),
\]
trong đó đồng cấu thứ nhất là đồng cấu chính tắc sinh bởi tích ngoài, còn đồng
cấu thứ hai được suy ra từ $s$. Khi ấy, kết luận suy ra từ (1.2) (i).

Kết hợp hai bổ đề trước, ta chứng minh được:

\result{BỔ ĐỀ}{2.6} Dưới các giả thiết sơ bộ của (2.1), giả sử với mọi
$\mathcal F$ ta có
\[
P_{\mathcal F}(n)=a_{\mathcal F}\frac{n^r}{r!}
+b_{\mathcal F}\frac{n^{r-1}}{(r-1)!}
+\text{các số hạng bậc $<r-1$},
\]
và các hệ số $a_{\mathcal F}$ bị chặn. Khi đó, các hệ số $b_{\mathcal F}$ bị
chặn dưới. Nếu các $b_{\mathcal F}$ bị chặn, thì các $\mathcal F_{(r)}$ tạo
thành một họ bị chặn.

Có thể giả sử các trường cơ sở $K$ của những bó $\mathcal F$ đều đóng đại số.
Trang bị cho mỗi $\operatorname{supp}\mathcal F_{(r)}$, là hợp của các thành
phần mang bậc $r$
của $\operatorname{supp}\mathcal F$, cấu trúc rút gọn cảm sinh. Khi đó, các
bậc của $\operatorname{supp}\mathcal F_{(r)}$ bị chặn trên bởi $a$; do đó,
theo (2.4), các $\operatorname{supp}\mathcal F_{(r)}$ tạo thành một tập hợp bị
chặn. Hơn nữa, với mỗi thành phần của
$\operatorname{supp}\mathcal F_{(r)}$, độ dài của $\mathcal F_{(r)}$ dọc theo
thành phần đó là $\leq a$; vì vậy, nếu $\mathcal J_{\mathcal F}$ là iđêan xác
định $\operatorname{supp}\mathcal F_{(r)}$, thì $\mathcal F_{(r)}$
\src{258}
có thể được xem như một môđun trên lược đồ con $Y_{\mathcal F}$ của $X$ được
xác định bởi $\mathcal J_{\mathcal F}^{a}$. Như trong bổ đề trước, ta cũng quy
về trường hợp $\mathcal F=\mathcal F_{(r)}$, do đó $\mathcal F$ xuất phát từ
một môđun trên $Y_{\mathcal F}$. Các $Y_{\mathcal F}$ tương ứng với một họ bị
chặn gồm các môđun thương của các $\mathcal O_{X_K}$, và vì thế xuất phát từ
một lược đồ con đóng $Y$ của một lược đồ $X\times_S T$. Khi đó, ta có thể áp
dụng (2.5) cho $Y/T$ và $\mathcal L\otimes_XY$, suy ra kết luận.

Bây giờ ta có thể chứng minh (2.2) bằng quy nạp theo cận trên $r$ của các số
$\dim\operatorname{supp}\mathcal F$. Mệnh đề là hiển nhiên với $r<0$; vì vậy,
giả sử $r\geq0$ và mệnh đề đã được chứng minh với mọi $r'<r$. Theo (2.6), các
$\mathcal F_{(r)}$ tạo thành một tập hợp bị chặn; do đó, theo (1.2) (i), các
hạt nhân của những đồng cấu
$\mathcal L_K\longrightarrow\mathcal F_{(r)}$ cũng tạo thành một tập hợp bị
chặn. Vì thế, tồn tại một môđun coherent $\mathcal L'$ trên $X$ sao cho các hạt
nhân đang xét, và do đó cả các
\[
N_r(\mathcal F)=\operatorname{Ker}(\mathcal F\longrightarrow\mathcal F_{(r)}),
\]
đều là những môđun thương của các môđun $\mathcal L'_K$. Vì các
$\mathcal F_{(r)}$ bị chặn, các $P_{\mathcal F_{(r)}}$ vẫn bị chặn, và công
thức
\[
P_{\mathcal F}=P_{\mathcal F_{(r)}}+P_{N_r}
\]
cho thấy các $P_{N_r}$ thỏa mãn cùng điều kiện $(b_s)$ như các
$P_{\mathcal F}$. Do đó, các $N_r$ thỏa mãn các điều kiện (a) và $(b_s)$, và
theo giả thiết quy nạp, các $(N_r)_{(s)}$ vẫn bị chặn. Mặt khác,
$\mathcal F_{(s)}$ là một mở rộng của $\mathcal F_{(r)}$ bởi
$(N_r)_{(s)}$; vì thế, theo (1.2) (ii), các $\mathcal F_{(s)}$ vẫn bị chặn.
Đối với khẳng định cuối cùng của (2.2), ta lưu ý rằng các hạt nhân $N_s$ của
$\mathcal F\longrightarrow\mathcal F_{(s)}$ vẫn bị chặn theo (1.2) (i), và hệ
số của số hạng chứa $n^{s-1}$ trong $P_{N_s}$ vẫn bị chặn; khi đó Bổ đề 2.6
chứng minh rằng hệ số của số hạng kế tiếp bị chặn dưới. Điều này kết thúc chứng
minh.

\fsection{3}{Các lược đồ Hilbert: định nghĩa, định lý tồn tại}

Cho $X$ là một tiền lược đồ trên một tiền lược đồ khác $S$, và $\mathcal F$ là
một môđun quasi-coherent trên $X$. Ký hiệu bởi
\[
\operatorname{Quot}(\mathcal F/X/S)
\]
tập hợp các môđun quasi-coherent thương của $\mathcal F$ phẳng trên $S$. Bây
giờ, cho $S'\longrightarrow S$ là một cấu xạ đổi cơ sở; đặt
$X'=X\times_S S'$ và
$\mathcal F'=\mathcal F\otimes_{\mathcal O_S}\mathcal O_{S'}$; do đó, $X'$
là một tiền lược đồ trên $S'$ được trang bị một môđun
\src{259}
quasi-coherent $\mathcal F'$, và ta có thể xét
$\operatorname{Quot}(\mathcal F'/X'/S')$. Ta đặt
\[
\underline{\operatorname{Quot}}_{\mathcal F/X/S}(S')
=\operatorname{Quot}(\mathcal F'/X'/S')
\qquad\text{(trong đó $X'=X\times_SS'$)}.
\]
Bây giờ, nếu $S''\longrightarrow S'$ là một $S$-cấu xạ, thì
$X''=X\times_SS''$ đẳng cấu với $X'\times_{S'}S''$, còn $\mathcal F''$ đẳng
cấu với $\mathcal F'\otimes_{\mathcal O_{S'}}\mathcal O_{S''}$. Vì hàm tử ảnh
ngược
\[
\mathcal G'\longmapsto
\mathcal G'\otimes_{\mathcal O_{S'}}\mathcal O_{S''}
\]
từ phạm trù các môđun quasi-coherent trên $X'$ vào phạm trù các môđun
quasi-coherent trên $X''$ là khớp phải và biến các môđun phẳng trên $S'$ thành
các môđun phẳng trên $S''$, ta thu được một ánh xạ tự nhiên
\[
\operatorname{Quot}(\mathcal F'/X'/S')\longrightarrow
\operatorname{Quot}(\mathcal F''/X''/S''),
\]
do đó $\underline{\operatorname{Quot}}_{\mathcal F/X/S}(S')$ là một hàm tử
phản biến theo $S'$ (một tiền lược đồ trên $S$), nhận giá trị trong phạm trù
các tập hợp. Từ đây về sau, ta giả sử $X$ xạ ảnh trên một $S$ Noether,
$\mathcal F$ coherent, và để đơn giản, ta chỉ xét các $S'$ trên $S$ là
Noether địa phương.

\result{ĐỊNH LÝ}{3.1} Dưới các điều kiện này, hàm tử phản biến
$\underline{\operatorname{Quot}}_{\mathcal F/X/S}$ trên phạm trù các
$S$-tiền lược đồ Noether địa phương được biểu diễn bởi một $S$-tiền lược đồ
$\operatorname{Quot}_{\mathcal F/X/S}$, là hợp rời nhau của một dãy các
$S$-lược đồ xạ ảnh (đặc biệt,
$\operatorname{Quot}_{\mathcal F/X/S}$ thuộc kiểu hữu hạn địa phương trên
$S$).

Ta sẽ thu được một phân rã như vậy theo cách sau. Cho $\mathcal O_X(1)$ là một
bó khả nghịch trên $X$, rất rộng tương đối với $S$. Với mỗi đa thức $P(n)$ có
các hệ số hữu tỉ, ký hiệu $\operatorname{Quot}^P(\mathcal F/X/S)$ là tập con
của $\operatorname{Quot}(\mathcal F/X/S)$ gồm các môđun thương coherent
$\mathcal G$ của $\mathcal F$, phẳng trên $S$, và có đa thức Hilbert tại mọi
$s\in S$ bằng $P$. Khi đó ta đặt
\[
\underline{\operatorname{Quot}}^P_{\mathcal F/X/S}(S')
=\operatorname{Quot}^P(\mathcal F'/X'/S')
\]
và như vậy thu được một hàm tử con của
$\underline{\operatorname{Quot}}_{\mathcal F/X/S}$. Tính bất biến của các đa
thức Hilbert được nhắc lại trong mục 2 kéo theo điều sau: để
$\underline{\operatorname{Quot}}_{\mathcal F/X/S}$ biểu diễn được, điều kiện
cần và đủ là các $\underline{\operatorname{Quot}}^P_{\mathcal F/X/S}$ đều
biểu diễn được; khi đó, $S$-tiền lược đồ
$\operatorname{Quot}_{\mathcal F/X/S}$ biểu diễn nó đẳng cấu với tiền lược đồ
là hợp rời nhau của các $\operatorname{Quot}^P_{\mathcal F/X/S}$ biểu diễn
những hàm tử $\underline{\operatorname{Quot}}^P_{\mathcal F/X/S}$. Vì vậy,
(3.1) là một hệ quả của định lý sau:
\src{260}
\result{ĐỊNH LÝ}{3.2} Với các ký hiệu trên, hàm tử
$\underline{\operatorname{Quot}}^P_{\mathcal F/X/S}$ được biểu diễn bởi một
$S$-tiền lược đồ xạ ảnh $\operatorname{Quot}^P_{\mathcal F/X/S}$.

Phần còn lại của mục này dành cho việc chứng minh (3.2).

Cho $\nu$ là một số nguyên. Với mọi $S'$ trên $S$, ta ký hiệu $A_\nu(S')$ là
tập hợp các môđun thương $\mathcal G=\mathcal F'/\mathcal K$ của
$\mathcal F'=\mathcal F\otimes_{\mathcal O_S}\mathcal O_{S'}$, coherent,
phẳng trên $S'$, và thỏa mãn các điều kiện sau:
\begin{align*}
\text{a.}\quad &R^if'_*(\mathcal G(n))=0 &&\text{với $i>0$ và $n\geq\nu$},\\
\text{b.}\quad &R^if'_*(\mathcal K(n))=0 &&\text{với $i>0$ và $n\geq\nu$},\\
\text{c.}\quad &f'_*(\mathcal K(\nu+k))
=\mathcal S'_k f'_*(\mathcal K(\nu)) &&\text{với $k\geq0$}.
\end{align*}
Đối với quan hệ cuối cùng này, ta giả sử $X$ được viết thành phổ nguyên tố
thuần nhất của một đại số phân bậc quasi-coherent $\mathcal S_*$ trên $S$, có
các bậc dương và được sinh bởi $\mathcal S_1$, rồi đặt
$\mathcal S'_* =\mathcal S_*\otimes_{\mathcal O_S}\mathcal O_{S'}$, sao cho
$X'$ là phổ nguyên tố thuần nhất của $\mathcal S'_*$. Để chứng minh (3.2), ta
dễ dàng quy về trường hợp $X=\mathbf P^r_S$ (vì $S$ là hợp của các tập mở $U$
sao cho $X|U$ là một lược đồ con đóng của $\mathbf P^r_U$ và
$\mathcal O_X(1)$ được cảm sinh bởi $\mathcal O_{\mathbf P^r_U}(1)$), và
trường hợp $\mathcal F$ có dạng $\mathcal O_{\mathbf P^r_S}^{N}$, do đó phẳng
trên $S$. Khi ấy, trong phần trên, các bó $\mathcal K$ cũng phẳng trên $S'$.
Từ các quan hệ Künneth ([5], III, \S~7), suy ra các điều kiện (a) và (b) ổn
định qua phép đổi cơ sở, đồng thời kéo theo rằng với $n\geq\nu$, việc tạo thành
$f'_*(\mathcal G(n))$ và $f'_*(\mathcal K(n))$ giao hoán với phép mở rộng cơ
sở (loc. cit.); do đó, (c) cũng ổn định qua phép mở rộng cơ sở.

Nói cách khác, $A_\nu(S')$ là một hàm tử phản biến theo $S'$, chính xác hơn là
một hàm tử con của
$A(S')=\underline{\operatorname{Quot}}^P_{\mathcal F/X/S}(S')$. Khi $\nu$
biến thiên, ta thu được một dãy tăng các tập con $A_\nu(S')$ của $A(S')$, có
hợp bằng $A(S')$ theo một định lý quen thuộc của Serre [5], III, \S~2. Bây giờ,
giả sử $\mathcal G$ là một môđun thương coherent của $\mathcal F'$, phẳng trên
$S'$, và $s$ là một điểm của $S'$ sao cho phép đổi cơ sở
$\operatorname{Spec}(k(s))\longrightarrow S'$ cho một môđun thương
$\mathcal G_s$ của $\mathcal F_s$ thỏa mãn các điều kiện (a), (b), (c), tức là
thuộc $A_\nu(\operatorname{Spec}(k(s)))$. Khi đó tồn tại một lân cận mở $U$
của $s$ sao cho các điều kiện ấy cũng được thỏa mãn đối với
$\mathcal G|f'^{-1}(U)$, tức là môđun thương này thuộc $A_\nu(U)$; đối với (a)
và (b), điều này suy ra từ «định lý về các hàm chỉnh hình» [5], III, \S~7, còn
(c) suy ra từ bổ đề Nakayama và từ việc trong mọi trường hợp ta biết rằng
$f'_*(\mathcal K(n+k))=\mathcal S'_k f'_*(\mathcal K(n))$
\src{261}
với $n$ đủ lớn và $k\geq0$ (xem [5], III, \S~2). Từ những nhận xét này, ta
suy ra điều sau (so sánh [4], IV). Để hàm tử $A$ biểu diễn được, điều kiện cần
và đủ là các hàm tử $A_\nu$ biểu diễn được; khi đó, $S$-tiền lược đồ $Q$ biểu
diễn $A$ là hợp tăng của các tiền lược đồ con mở $Q_\nu$ biểu diễn các
$A_\nu$.

Cho
\[
M_* =\bigoplus_{n\geq0}f_*(\mathcal F(n))=\mathcal S_*^{N},
\]
do đó
\[
M'_* =M_*\otimes_{\mathcal O_S}\mathcal O_{S'}
=\bigoplus_{n\geq0}f'_*(\mathcal F'(n))=\mathcal S'_*{}^{N}.
\]
Từ (a) suy ra:

\noindent (a') $f'_*(\mathcal G(n))$ tự do địa phương hạng $P(n)$ với
$n\geq\nu$, [5], III, \S~7, và từ (b), với $i=1$, suy ra:

\noindent (a'') $f'_*(\mathcal G(n))$ là một môđun thương của $M'_n$.

Hơn nữa, việc biết môđun thương này khi $n=\nu$, theo (c), xác định các môđun
con $f'_*(\mathcal K(n))$ của $M'_n$ với $n\geq\nu$, do đó xác định
$\mathcal K$ và tiếp theo là $\mathcal G$. Như vậy, ta thu được một ánh xạ đơn
\[
A_\nu(S')\longrightarrow\operatorname{Grass}_{P(\nu)}(M'_\nu)
\]
từ $A_\nu(S')$ vào tập hợp các môđun thương tự do địa phương hạng $P(\nu)$ của
$M'_\nu$, và do đó một đồng cấu theo hàm tử
\[
i_\nu:A_\nu(S')\longrightarrow
\underline{\operatorname{Grass}}_{P(\nu)}(M_\nu)(S'),
\]
trong đó hàm tử ở vế phải được biểu diễn bởi lược đồ Grassmann
$\operatorname{Grass}_{P(\nu)}(M_\nu)$ (so sánh [4], V), là một lược đồ xạ
ảnh trên $S$. Ta khẳng định:

\result{BỔ ĐỀ}{3.3} $A_\nu(S')$ là một hàm tử biểu diễn được, và cấu xạ
$Q_\nu\longrightarrow\operatorname{Grass}_{P(\nu)}(M_\nu)$ biểu diễn đồng cấu
$i_\nu$ là một phép nhúng (điều này kéo theo $Q_\nu$ tựa xạ ảnh trên $S$).

Khẳng định này tương đương với khẳng định sau (so sánh [4], IV): giả sử cho
một môđun thương $N$ của $M'_\nu$, tự do địa phương hạng $P(\nu)$; khi đó tồn
tại một tiền lược đồ con $Z$ của $S'$ sao cho với mọi tiền lược đồ Noether địa
phương $T'$
\src{262}
trên $S'$, ảnh ngược của $N$ trên $T'$ thuộc
$\operatorname{Im}A_\nu(T')$ khi và chỉ khi $T'\longrightarrow S'$ phân tích
qua tiền lược đồ con $Z$. Thay đổi ký hiệu, ta có thể giả sử $S'=S$, tức là ta
có một môđun thương $N_\nu$ của $M_\nu$ theo một môđun con $R_\nu$. Để môđun
thương này phát sinh từ một phần tử của $A_\nu(S)$, điều kiện cần và đủ là nó
thỏa mãn hai điều kiện sau:

\medskip
\noindent (i) $M_{\nu+k}/\mathcal S_kR_\nu$ tự do địa phương hạng
$P(\nu+k)$ với $k\geq0$.

\smallskip
\noindent (ii) Bó con $\mathcal K$ của $\mathcal F$, được xác định bởi môđun
con phân bậc
\[
R_* =\sum_{k\geq0}\mathcal S_kR_\nu
\]
của $M_*$ ([5], II, \S~3), và môđun thương
$\mathcal G=\mathcal F/\mathcal K$ thỏa mãn các điều kiện (a) và (b) ở trên.

Các điều kiện này hiển nhiên là cần thiết; ngược lại, nếu chúng được thỏa mãn,
thì bó $\mathcal G$ xác định trong (ii), do đẳng cấu với bó liên kết với
$\mathcal S_*$-môđun phân bậc $N_*$ là tổng của các
$M_{\nu+k}/\mathcal S_kR_\nu$, phẳng trên $S$ (vì các thớ của nó là những hạng
tử trực tiếp của các môđun thu được bằng cách địa phương hóa $N$ tại các
iđêan nguyên tố thuần nhất của $\mathcal S_*$), và có đa thức Hilbert $P$ theo
(i). Xét đến (ii), ta
thấy khi ấy (a') và (a'') được thỏa mãn, do đó, với $n\geq\nu$, đồng cấu tự
nhiên $N_n\longrightarrow f_*(\mathcal G(n))$ là một đồng cấu toàn ánh giữa
các môđun tự do địa phương có cùng hạng, nên là một đẳng cấu; vì vậy
\[
f_*(\mathcal K(\nu+k))=\mathcal S_kR_\nu
\quad\text{với mọi $k\geq0$},
\]
điều này chứng minh rằng $\mathcal G\in A_\nu(S)$ và $R_\nu$ là phần tử của
$\operatorname{Grass}_{P(\nu)}(M_\nu)$ được xác định bởi $\mathcal G$.

Tiêu chuẩn (i)--(ii) này cũng áp dụng cho tình huống thu được sau một phép đổi
cơ sở $S'\longrightarrow S$. Trước hết, ta sẽ chứng minh rằng việc điều kiện
(i) được thỏa mãn sau phép đổi cơ sở $S'\longrightarrow S$ có thể biểu thị
bằng điều kiện rằng $S'\longrightarrow S$ phân tích qua một tiền lược đồ con
$Z$ nào đó của $S$; sau khi có kết quả này, ta quy về (bằng cách thay $S$ bởi
$Z$) trường hợp điều kiện (i) đã được thỏa mãn trên $S$, và vì nó ổn định qua
phép đổi cơ sở nên chỉ còn phải biểu thị điều kiện (ii). Nhưng khi ấy, nếu $U$
ký hiệu tập hợp các $s\in S$ sao cho đối đồng điều của các bó cảm sinh trên thớ
$X_s$ bởi $\mathcal G(n)$ và $\mathcal K(n)$ triệt tiêu ở bậc $>0$ với
$n\geq\nu$, ta đã lưu ý rằng $U$ là mở, và điều kiện (ii) sẽ được thỏa mãn sau
một phép đổi cơ sở $S'\longrightarrow S$ khi và chỉ khi
$S'\longrightarrow S$ phân tích qua $U$; điều này chứng minh Bổ đề 3.3. Do đó,
còn phải chứng minh bổ đề sau:

\result{BỔ ĐỀ}{3.4} Cho $S$ là một tiền lược đồ Noether địa phương, được trang
bị một đại số phân bậc quasi-coherent $\mathcal S_*$ có các bậc dương, được
sinh bởi $\mathcal S_1$; cho $M_*$ là một $\mathcal S_*$-môđun phân bậc quasi-coherent
hữu hạn sinh, $P$ là một đa thức có hệ số hữu tỉ, và $\nu$ là một số nguyên.
Khi đó tồn tại một tiền lược đồ con $Z$ của $S$ (hiển nhiên là duy nhất)
\src{263}
có tính chất sau: với mọi tiền lược đồ $S'$ trên $S$, điều kiện cần và đủ để
$M_n\otimes_{\mathcal O_S}\mathcal O_{S'}$ tự do địa phương hạng $P(n)$ với
mọi $n\geq\nu$ là $S'\longrightarrow S$ phân tích qua $Z$.

Ta hiển nhiên có thể giả sử $S$ affine, do đó là Noether. Khi ấy:

\result{BỔ ĐỀ}{3.5} Với mọi số nguyên $N\geq\nu$, đặt $U_N$ là tập mở của $S$
gồm các $s\in S$ sao cho
\[
\operatorname{rank}_{k(s)}
M_{ns}\otimes_{\mathcal O_{S,s}}k(s)\leq P(n)
\quad\text{với mọi $\nu\leq n\leq N$}.
\]
Khi đó dãy giảm các tập mở $U_N$ là dừng.

Ta biết [3], IV, rằng $S$ có một phân hoạch hữu hạn thành các tiền lược đồ con
rút gọn $S_i$, sao cho mỗi $M\otimes_{\mathcal O_S}\mathcal O_{S_i}$ phẳng
trên $S_i$. Do đó, ta có thể giả sử $M$ phẳng, nên các $M_n$ phẳng. Cuối cùng,
ta hiển nhiên có thể giả sử $S$ liên thông. Nhưng khi ấy tồn tại một số nguyên
$n_0$ và một đa thức $Q$ sao cho
\[
\operatorname{rank}_{k(s)}
M_{ns}\otimes_{\mathcal O_{S,s}}k(s)=Q(n)
\quad\text{với $n\geq n_0$},
\]
[5], III, \S~7. Trước hết giả sử $P\neq Q$; khi đó
$P(n)\neq Q(n)$ với $n$ lớn, nên hiển nhiên $U_N=\varnothing$ với $N$ lớn;
huống chi dãy các $U_N$ là dừng. Ngược lại, nếu $P=Q$, thì
$U_N=U_{n_0}$ với $N\geq n_0$, và dãy các $U_N$ vẫn là dừng.

Đặc biệt, tập $U_\infty$ gồm các $s\in S$ sao cho
\[
\tag{*}
\operatorname{rank}_{k(s)}
M_{ns}\otimes_{\mathcal O_{S,s}}k(s)\leq P(n)
\quad\text{với mọi $n\geq\nu$},
\]
là mở vì nó là giao của các $U_N$. Khi ấy, để chứng minh (3.4), ta có thể thay
$S$ bởi tập mở $U$, nhờ đó quy về trường hợp bất đẳng thức (*) được thỏa mãn
tại mọi $s\in S$. Hơn nữa:

\result{BỔ ĐỀ}{3.6} Cho $M$ là một môđun trên một tiền lược đồ Noether địa
phương $S$, và $r$ là một số nguyên. Khi đó tồn tại một tiền lược đồ con $Z$
của $S$ (hiển nhiên là duy nhất) có tính chất sau: với mọi $S'$ trên $S$, điều
kiện cần và đủ để $M\otimes_{\mathcal O_S}\mathcal O_{S'}$ tự do địa phương
hạng $r$ là $S'\longrightarrow S$ phân tích qua $Z$. Nếu
\[
\operatorname{rank}_{k(s)}
M_s\otimes_{\mathcal O_{S,s}}k(s)\leq r
\quad\text{với mọi $s$},
\]
thì $Z$ là một tiền lược đồ con đóng của $S$ (giả sử $M$ coherent).
\src{264}
Thật vậy, lập luận trên quy ta về trường hợp bất đẳng thức (*) đúng với mọi
$s\in S$ (nếu cần, thay $S$ bằng phần mở gồm các $s$ tại đó bất đẳng thức
được thỏa mãn). Khi ấy ta có thể giả sử $M$ nằm trong một dãy khớp
\[
\mathcal O_S^q\longrightarrow\mathcal O_S^r\longrightarrow M\longrightarrow0,
\]
và điều kiện đang xét đối với các $S'$ trên $S$ cũng có nghĩa là, trong dãy
khớp tương ứng $\mathcal O_{S'}^q\longrightarrow\mathcal O_{S'}^r
\longrightarrow M'\longrightarrow0$, mũi tên thứ hai là một đẳng cấu, tức là
mũi tên thứ nhất bằng không. Khi đó, tiền lược đồ con đóng $Z$ của $S$ được
xác định bởi iđêan sinh bởi các hệ số của ma trận xác định đồng cấu
$\mathcal O_S^q\longrightarrow\mathcal O_S^r$ thỏa mãn điều kiện mong muốn.

Trở lại chứng minh (3.4) tại chỗ ta đã dừng, gọi $Z_n$ là lược đồ con đóng
của $S$ được liên kết, theo (3.6), với môđun $M_n$ và số nguyên $r=P(n)$, và
gọi $Z'_N$ là $\operatorname{Inf}$ của các $Z_n$ với
$\nu\leq n\leq N$. Khi ấy các $Z'_N$ tạo thành một dãy giảm các lược đồ con
đóng của $Z$, nên tất yếu dừng. Gọi $Z$ là giá trị không đổi của các
$Z'_N$ khi $N$ lớn. Đó là $Z$ cần tìm trong (3.4). Điều này kết thúc chứng
minh (3.4), và do đó cả chứng minh (3.3).

Vậy ta đã chứng minh rằng
$\underline{\operatorname{Quot}}^P_{\mathcal F/X/S}$ được biểu diễn bởi một
$S$-tiền lược đồ $Q$ là hợp tăng của các tiền lược đồ con mở $Q_\nu$ tựa xạ
ảnh trên $S$. Để tiến thêm, phải viện dẫn Định lý 2.1, từ đó dễ dàng suy ra
$Q$ là quasi-compact (vì nó là ảnh của tiền lược đồ $S'$ thuộc kiểu hữu hạn
trên $S$ tham số hóa họ các bó thương của $\mathcal F_K$ có đa thức Hilbert
$P$). Do đó $Q$ bằng một trong các $Q_\nu$, vì thế tựa xạ ảnh trên $S$. Để
chứng minh nó xạ ảnh trên $S$, chỉ còn phải chứng minh nó riêng trên $S$; để
làm điều này, chỉ cần viện dẫn tiêu chuẩn định giá cho tính riêng dưới dạng
[5], II, 7.3.8. Chỉ cần kiểm tra điều sau:

\result{BỔ ĐỀ}{3.7} Cho $S$ là phổ của một vành định giá rời rạc, $s$ là điểm
tổng quát của nó, $X$ là một tiền lược đồ trên $S$, $\mathcal F$ là một môđun
quasi-coherent trên $X$, và $\mathcal G_s$ là một môđun thương quasi-coherent
của $\mathcal F_s=\mathcal F\otimes_{\mathcal O_S}k(s)$ trên $X_s$. Khi đó
tồn tại duy nhất một môđun thương quasi-coherent $\mathcal G$ của
$\mathcal F$, phẳng trên $S$, sao cho hạn chế của nó lên $X_s$ là
$\mathcal G_s$.
\src{265}
Thật vậy, nếu $\mathcal G_s=\mathcal F_s/\mathcal K_s$, chỉ cần xét bó con lớn
nhất $\mathcal K$ của $\mathcal F$ cảm sinh $\mathcal K_s$, [5], I, 9.4.2, và
lấy $\mathcal G=\mathcal F/\mathcal K$. Dễ dàng kiểm tra rằng bó này thỏa mãn
yêu cầu.

Định lý 3.2, và do đó (3.1), được chứng minh hoàn toàn.

Đồng thời, chứng minh cũng cho ta điều sau:

\result{MỆNH ĐỀ}{3.8} Trong các điều kiện của (3.2), đặt
$Q=\operatorname{Quot}^P_{\mathcal F/X/S}$, $X_Q=X\times_SQ$,
$\mathcal F_Q=\mathcal F\otimes_{\mathcal O_S}\mathcal O_Q$, và cho
$\mathcal G$ là môđun thương coherent của $\mathcal F_Q$, phẳng trên $Q$, có
đa thức Hilbert tương đối $P$, sao cho $(Q,\mathcal G)$ biểu diễn hàm tử
$\underline{\operatorname{Quot}}^P_{\mathcal F/X/S}$. Tồn tại một số nguyên
$\nu$ sao cho, với $n\geq\nu$, $(f_Q)_*(\mathcal G(n))$ là một môđun trên $Q$
tự do địa phương hạng $P(n)$ và rất rộng tương đối với $S$, tức là xác định
một phép nhúng của $Q$ vào một lược đồ Grassmann
$\operatorname{Grass}_{P(n)}(M)$ trên $S$. Đặc biệt, với $n\geq\nu$, bó
\[
\bigwedge^{P(n)}(f_Q)_*(\mathcal G(n))
\]
trên $Q$ là khả nghịch và rất rộng tương đối với $S$.

Thật vậy, như đối với (3.2), ta có thể quy về trường hợp $\mathcal F$ phẳng
trên $S$, và khi đó chỉ cần lấy một số nguyên $\nu$ sao cho $A_\nu=A$ (với các
ký hiệu trước).

Trường hợp áp dụng quan trọng nhất của (3.2) là khi
$\mathcal F=\mathcal O_X$. Khi ấy ta viết
\[
\operatorname{Quot}_{\mathcal O_X/X/S}=\operatorname{Hilb}_{X/S},
\qquad
\operatorname{Quot}^P_{\mathcal O_X/X/S}=\operatorname{Hilb}^P_{X/S},
\]
do đó có một phân rã
\[
\operatorname{Hilb}_{X/S}=\coprod_P\operatorname{Hilb}^P_{X/S}.
\]
Theo định nghĩa, $\operatorname{Hilb}_{X/S}$ biểu diễn hàm tử
$\underline{\operatorname{Hilb}}_{X/S}(S')$, là tập hợp các tiền lược đồ con
đóng của $X'=X\times_SS'$ phẳng trên $S'$, và
$\operatorname{Hilb}^P_{X/S}$ biểu diễn hàm tử con tương ứng với các tiền lược
đồ con đóng có một đa thức Hilbert cho trước $P$. Các tiền lược đồ này lần lượt
còn được gọi là \emph{tiền lược đồ Hilbert} của $X$ trên $S$ và
\emph{tiền lược đồ Hilbert chỉ số $P$}. Thuật ngữ này được biện minh bởi vai
trò của các đa thức Hilbert trong lý thuyết. Về bản chất, sự khác biệt của
chúng với các đa tạp Chow cổ điển (dùng để tham số hóa các chu trình, chứ không
phải các đa tạp) tương tự như sự khác biệt giữa
\src{266}
vành Chow của các lớp chu trình trên một đa tạp và vành các lớp bó trên đa
tạp đó (như vành này xuất hiện trong định lý Riemann--Roch [1]); thật vậy,
lưu ý rằng khi $X=\mathbf P^r_S$, với $S$ là phổ của một trường, việc biết đa
thức Hilbert của một môđun coherent $\mathcal F$ trên $X$ cũng tương đương với
việc biết các lớp Chern của $\mathcal F$, hay lớp của $\mathcal F$ trong vành
các lớp bó coherent trên $X$.

\result{CÁC NHẬN XÉT}{3.9} Cũng cần lưu ý rằng việc xây dựng các
$\operatorname{Quot}_{\mathcal F/X/S}$ và
$\operatorname{Quot}^P_{\mathcal F/X/S}$ đã được quy về trường hợp
$X=\mathbf P^r_S$, $\mathcal F=\mathcal O_X^N$, trong đó
$\mathcal O_X(1)$ là bó rất rộng thông thường; chính xác hơn, các
$\operatorname{Quot}_{\mathcal F/X/S}$ tổng quát là các tiền lược đồ con đóng
của những tiền lược đồ vừa nêu. Vì việc tạo thành các
$\operatorname{Quot}^P_{\mathcal F/X/S}$ hiển nhiên tương thích với các phép
đổi cơ sở $S'\longrightarrow S$, ta thấy có thể quy tiếp về trường hợp
$S=\operatorname{Spec}(\mathbf Z)$. Sau cùng, ta được quy về việc nghiên
cứu các lược đồ xạ ảnh sau đây trên $\mathbf Z$:
\[
\mathfrak Q^P_{r,N}=
\operatorname{Quot}^{P}_{(\mathcal O_{\mathbf P^r_{\mathbf Z}})^N/
\mathbf P^r_{\mathbf Z}/\operatorname{Spec}(\mathbf Z)},
\]
và đặc biệt hơn là các lược đồ Hilbert tuyệt đối:
\[
\operatorname{Hilb}^P_r=\mathfrak Q^P_{r,1}.
\]
Sẽ rất đáng quan tâm nếu nghiên cứu chi tiết hơn các lược đồ này, bắt đầu bằng
việc xác định các thành phần liên thông của chúng (các lược đồ này có liên
thông không?), và các thành phần bất khả quy của
chúng (theo Serre [7], có thể có những thành phần bất khả quy nằm hoàn toàn
phía trên một số nguyên tố $p\neq0$). Nhắc lại câu hỏi của Weil: liệu các thành
phần bất khả quy của các thớ của $\operatorname{Hilb}^P_r$ trên các
$s\in\operatorname{Spec}(\mathbf Z)$ có tương ứng với các mở rộng
``chính quy'' của trường nguyên tố hay không, tức là liệu chúng có ``liên thông
tương đối'' hay không. Không loại trừ khả năng những câu hỏi này dễ tiếp cận
đối với các lược đồ Hilbert hơn là đối với các ``đa tạp Chow''.

\fsection{4}{Các biến thể}

\noindent a. Trong các điều kiện của (3.1), cho $U$ là một tập mở trong $X$,
và ký hiệu $A'$ là hàm tử con của
$A=\underline{\operatorname{Quot}}_{\mathcal F/X/S}$ sao cho $A'(S')$ là tập
các môđun thương $\mathcal G$ của $\mathcal F'$, phẳng trên $S'$, có giá chứa
trong $U'$. Ta thấy ngay rằng $A'$ được biểu diễn bởi một tập mở của tiền lược
đồ $\operatorname{Quot}_{\mathcal F/X/S}$ biểu diễn $A$. Do đó các định lý
(3.1) và (3.2)
\src{267}
vẫn còn đúng khi giả sử $X$ tựa xạ ảnh trên $S$ thay vì xạ ảnh trên $S$, đồng
thời trong các kết luận cũng thay từ ``xạ ảnh'' bằng ``tựa xạ ảnh'', và giờ
đây ký hiệu
$\underline{\operatorname{Quot}}_{\mathcal F/X/S}(S')$ là tập các môđun
thương coherent $\mathcal G$ của $\mathcal F'$, phẳng trên $S'$, có giá riêng
trên $S'$.

\medskip
\noindent b. Nói chung, có thể áp đặt lên các môđun thương $\mathcal G$ của
$\mathcal F'=\mathcal F\otimes_{\mathcal O_S}\mathcal O_{S'}$, phẳng trên
$S'$, mọi loại điều kiện tự nhiên bổ sung ổn định qua phép đổi cơ sở, qua đó thu
được các hàm tử con của
$\underline{\operatorname{Quot}}_{\mathcal F/X/S}$ mà ta muốn biểu diễn. Trong
nhiều trường hợp, tiêu chuẩn thông thường cho phép chứng minh rằng các hàm tử
thu được vẫn được biểu diễn bởi các tập mở của
$\operatorname{Quot}_{\mathcal F/X/S}$. Đặc biệt, điều này đúng khi áp đặt một
trong các tính chất bổ sung sau:

\smallskip
\noindent $1^\circ$ Chiều của mỗi chu trình nguyên tố liên kết với các môđun
$\mathcal G_{s'}$ ($s'\in S'$), được cảm sinh trên các thớ $X'_{s'}$, thuộc
một tập số nguyên đã cho.

\smallskip
\noindent $2^\circ$ (Khi $\mathcal F=\mathcal O_X$, do đó $\mathcal G$ tương
ứng với một tiền lược đồ con đóng $Y$ của $X'$); $Y$ là một tiền lược đồ trơn
([3], IV) trên $S$, hoặc chuẩn tắc trên $S$ (tức là các thớ $Y_{s'}$ chuẩn
tắc ``trên $k(s')$'', nghĩa là vẫn chuẩn tắc sau mọi mở rộng của trường cơ
sở), hoặc (khi $X$ phẳng trên $S$) là một giao đầy đủ địa phương theo cơ sở
$S$ trong $X$, xét tương đối trên $S$ (tức là các thớ $Y_{s'}$ là các giao
đầy đủ địa phương trong các $X_{s'}$).

Những điều kiện khác sẽ liên quan đến các tính chất đối đồng điều
của các môđun $\mathcal G_{s'}$ được cảm sinh trên các $X'_{s'}$, v.v. Dĩ
nhiên, hội của các điều kiện mà mỗi điều kiện được biểu diễn bởi một tập mở
$U_i$ của $\operatorname{Quot}_{\mathcal F/X/S}$ được biểu diễn bởi tập mở thu
được bằng cách lấy giao các tập ấy. Chẳng hạn, với mỗi $S'$ trên $S$, xét tập
các tiền lược đồ con đóng $Y$ của $X'=X\times_SS'$ mà mỗi tiền lược đồ là một
phủ étale ([3], I) hạng cho trước $r$ của $S'$, ta thu được một hàm tử phản biến
biểu diễn được theo biến $S'$.

\medskip
\noindent c. Các tiền lược đồ $\operatorname{Hom}_S(X,Y)$,
$\Pi_{X/S}(Z/X)$, $\operatorname{Isom}_S(X,Y)$, được định nghĩa trong [2], II,
C, n\textsuperscript{o} 2, tồn tại dưới các giả thiết xạ ảnh thích hợp và được
hiện thực hóa dưới dạng các tập mở trong những tiền lược đồ Hilbert thích hợp.
Vì ta có
\[
\operatorname{Hom}_S(X,Y)=\Pi_{X/S}((X\times_SY)/X),
\]
trường hợp $\operatorname{Hom}_S(X,Y)$ được quy về trường hợp
$\Pi_{X/S}(Z/X)$. Khi đó, lưu ý rằng với mỗi $S'$ trên $S$, tập các tiết diện
của $Z'=Z\times_SS'$ trên $X'=X\times_SS'$ tương ứng song ánh với tập các tiền
lược đồ con $\Gamma$ của $Z'$ (nhất thiết đóng nếu $Z$ tách trên $X$) sao cho
cấu xạ $\Gamma\longrightarrow X'$ cảm sinh bởi $Z'\longrightarrow X'$ là một
đẳng cấu. Bằng cách này,
\src{268}
khi $X$ phẳng và riêng trên $S$, còn $Z$ tựa xạ ảnh trên $S$,
$\Pi_{X/S}(Z/X)$ tồn tại và được hiện thực hóa dưới dạng một tiền lược đồ con mở của
$\operatorname{Hilb}_{Z/S}$. Do đó, khi $X$ xạ ảnh và phẳng trên $S$, còn $Y$
tựa xạ ảnh trên $S$, $\operatorname{Hom}_S(X,Y)$ tồn tại và được hiện thực hóa
dưới dạng một tiền lược đồ con mở của
$\operatorname{Hilb}_{(X\times_SY)/S}$. Khi $X$ và $Y$ đều xạ ảnh trên $S$,
suy ra ngay $\operatorname{Isom}_S(X,Y)$ cũng tồn tại và được biểu diễn bởi
một tập mở của $\operatorname{Hom}_S(X,Y)$. Tương tự, khi $X$ phẳng và xạ ảnh
trên $S$, còn $Y$ tựa xạ ảnh trên $S$, $S$-tiền lược đồ
$\operatorname{Imm}_S(X/Y)$, tương ứng với hàm tử con của hàm tử được
$\operatorname{Hom}_S(X,Y)$ biểu diễn, gồm các $S'$-cấu xạ
$X'\longrightarrow Y'$ là các phép nhúng, cũng được biểu diễn bởi một tập mở
của $\operatorname{Hom}_S(X,Y)$.

Cho $\mathcal L$ (tương ứng $\mathcal M$) là một bó khả nghịch trên $X$ (tương
ứng trên $Y$), rất rộng tương đối với $S$; từ đó có một bó rất rộng
$\mathcal L\otimes_{\mathcal O_S}\mathcal M$ trên $X\times_SY$ tương đối với
$S$. Vì vậy, với mọi đa thức $P$ có hệ số hữu tỉ,
$\operatorname{Hilb}^P_{(X\times_SY)/S}$ được xác định và là một tiền lược đồ
tựa xạ ảnh trên $S$. Do đó nó cảm sinh trên $\operatorname{Hom}_S(X,Y)$ một
tập vừa mở vừa đóng, tựa xạ ảnh trên $S$, mà ta ký hiệu là
$\operatorname{Hom}_S(X,Y)^P$. Vậy các tiết diện của
$\operatorname{Hom}_S(X,Y)^P$ trên $S$ là các $S$-cấu xạ
$g:X\longrightarrow Y$ sao cho với mọi số nguyên $n$, ta có
\[
\chi\!\left(\left(\mathcal L\otimes_{\mathcal O_X}g^*(\mathcal M)\right)^{
\otimes n}\right)=P(n).
\]

Theo cách đó, ta thu được những tổng quát hóa của định lý Matsusaka, khẳng định
rằng các tự đẳng cấu của một đa tạp xạ ảnh ``phân cực'' tạo thành một nhóm đại
số. Ở đây, khẳng định này hiển nhiên mang một ý nghĩa chính xác hơn, vì ta có
một định nghĩa của nhóm ấy như nghiệm của một bài toán phổ dụng. Cũng cần lưu ý
rằng, trên một trường đóng đại số, nhóm các tự đẳng cấu được xét trước đây là
nhóm thu được từ nhóm ``thật'' được định nghĩa ở đây bằng cách lấy thương theo
iđêan các phần tử lũy linh; điều này giải thích vì sao khó có thể thực hiện các
phép dựng cũ trên một trường cơ sở không hoàn hảo: iđêan các phần
tử lũy linh xuất hiện sau một phép mở rộng trường cơ sở không nhất thiết
``được xác định trên $k$''. Nhận xét này cũng áp dụng cho phần lớn các phép dựng
theo lối cũ.

\fsection{5}{Nghiên cứu vi phân các lược đồ Hilbert}

Nghiên cứu này suy ra từ kết quả sau:
\src{269}
\result{MỆNH ĐỀ}{5.1} Cho $S$ là một tiền lược đồ, $S_0$ là một tiền lược đồ con
được xác định bởi một iđêan quasi-coherent $\mathcal J$ có bình phương bằng
không, $X$ là một $S$-tiền lược đồ, $\mathcal F$ là một môđun quasi-coherent
trên $X$, đặt $X_0=X\times_SS_0$ và
$\mathcal F_0=\mathcal F\otimes_{\mathcal O_S}\mathcal O_{S_0}$, và cuối cùng
cho $\mathcal G_0=\mathcal F_0/\mathcal K_0$ là một môđun thương
quasi-coherent của $\mathcal F_0$, phẳng trên $S_0$. Với mỗi tập mở $U$ của
$X$, gọi $\mathcal E(U)$ là tập các môđun thương quasi-coherent $\mathcal G$
của $\mathcal F|U$, phẳng trên $S$, và thỏa mãn
$\mathcal G\otimes_{\mathcal O_S}\mathcal O_{S_0}=\mathcal G_0$; do đó, khi
$U$ thay đổi, các $\mathcal E(U)$ là các tiết diện của một bó $\mathcal E$
trên $U$. Khi ấy, bó các nhóm
\[
\mathcal A=\operatorname{Hom}_{\mathcal O_{X_0}}
(\mathcal K_0,\mathcal G_0\otimes_{\mathcal O_{S_0}}\mathcal J)
\]
tác động tự nhiên lên $\mathcal E$, nhờ đó bó này trở thành một bó
``thuần nhất chính một cách hình thức dưới $\mathcal A$'' (nghĩa là, với mọi
tập mở $U$ trong $X$, $\mathcal E(U)$ hoặc là rỗng, hoặc là một tập thuần nhất
chính dưới $\mathcal A(U)$).

Từ đó suy ra:

\result{HỆ QUẢ}{5.2} Giả sử địa phương trên $X$ tồn tại một phép nâng
$\mathcal G$ của $\mathcal G_0$ thành một môđun thương của $\mathcal F$ phẳng
trên $S$ (nghĩa là các thớ của bó $\mathcal E$ không rỗng). Khi đó tồn tại một
lớp cản trở chính tắc
\[
c(\mathcal G_0)\in H^1(X,\mathcal A),
\]
mà sự triệt tiêu của nó là điều kiện cần và đủ để tồn tại một phép nâng toàn
cục $\mathcal G$ của $\mathcal G_0$ thành một môđun thương của $\mathcal F$
phẳng trên $S$. Nếu lớp này bằng không, thì tập $\mathcal E(X)$ gồm mọi phép
nâng khả dĩ là một tập thuần nhất chính dưới
\[
\mathcal A(X)=\operatorname{Hom}_{\mathcal O_X}
(\mathcal K_0,\mathcal G_0\otimes_{\mathcal O_{S_0}}\mathcal J).
\]

Đặc biệt, sự tồn tại của một phép nâng toàn cục được bảo đảm nếu
$H^1(X,\mathcal A)=0$.

\result{HỆ QUẢ}{3.3} Giả sử
$Q=\operatorname{Quot}_{\mathcal F/X/S}$ tồn tại (xem 4 (a)) -- chẳng hạn,
$X$ tựa xạ ảnh trên tiền lược đồ Noether địa phương $S$ và $\mathcal F$ coherent. Cho
$x\in Q$ tương ứng với một mở rộng trường thặng dư $K=k(x)$ của $k(s)$
($s\in S$):
\src{270}
do đó, $x$ được xác định bởi một môđun thương coherent
$\mathcal G_0=\mathcal F_0/\mathcal K_0$ của môđun
$\mathcal F_0=\mathcal F_K$ trên $K$-tiền lược đồ $X_K$. Gọi $\mathcal A$ là
bó coherent trên $X_K$ được xác định bởi
\[
\mathcal A=\operatorname{Hom}_{\mathcal O_{X_0}}(\mathcal K_0,\mathcal G_0).
\]
Khi đó, không gian tiếp tuyến Zariski của thớ $Q_s$ tại điểm $x$ (đối ngẫu trên
$K$ của $\mathfrak m/\mathfrak m^2$, trong đó $\mathfrak m$ là iđêan cực đại
của $\mathcal O_{Q_K,x}$) đẳng cấu chính tắc với $H^0(X_K,\mathcal A)$.

Mô tả không gian tiếp tuyến Zariski được tổng quát hóa như sau. Với một
$S$-cấu xạ cho trước $g:S'\longrightarrow Q$, tức là một tiết diện $g'$ của
$Q'=Q\times_SS'$ trên $S'$, môđun
\[
\Omega=g^*(\Omega^1_{Q/S})=g'^*(\mathcal J/\mathcal J^2)
\]
(trong đó $\mathcal J$ là iđêan trên $Q'$ được xác định bởi tiết diện $g'$ của
$Q'$ trên $S'$) được đặc trưng bởi công thức sau, có tính hàm tử theo môđun coherent
$\mathcal M$ trên $S'$:
\[
\operatorname{Hom}_{\mathcal O_{S'}}(\Omega,\mathcal M)
\simeq H^0(X',\mathcal A),
\]
trong đó $\mathcal A$ vẫn là môđun trên $X'=X\times_SS'$ được xác định bởi
\[
\mathcal A=\operatorname{Hom}_{\mathcal O_{X'}}
(\mathcal K,\mathcal G\otimes_{\mathcal O_S}\mathcal M),
\]

($\mathcal G=\mathcal F'/\mathcal K$ là môđun thương của
$\mathcal F'=\mathcal F\otimes_{\mathcal O_S}\mathcal O_{S'}$ tương ứng với
$g$). Thật vậy, chỉ cần áp dụng (5.1), trong đó thay $S_0$ bởi $S'$, thay $S$
bởi tiền lược đồ $D(\mathcal M)=(S',\mathcal O_{S'}+\mathcal M)$, với
$\mathcal M$ được xem như một iđêan có bình phương bằng không.

Khi $\mathcal F=\mathcal O_X$ trong (5.1), dữ liệu của
$\mathcal G_0$ tương ứng với một tiền lược đồ con đóng $Y_0$ của
$X_0$, phẳng trên $S_0$ và được xác định bởi iđêan
$\mathcal J_0=\mathcal K_0$; khi đó (*) cho
\[
\mathcal A=\operatorname{Hom}_{\mathcal O_{X_0}}
(\mathcal J_0/\mathcal J_0^2,
\mathcal O_{Y_0}\otimes_{\mathcal O_{S_0}}\mathcal J),
\]
trong đó $\mathcal J/\mathcal J^2$ được hiểu là bó đối pháp tuyến của $Y_0$
trong $X_0$, cũng được ký hiệu là
\src{271}
$\check{\mathcal N}_{Y_0/X_0}$; khi đó nên xem $\mathcal A$ như một môđun
trên $Y_0$ và tính $H^0$ và $H^1$ trên $Y_0$. Hơn nữa, nếu $Y_0$ là một giao
đầy đủ địa phương trong $X_0$, với $X$ phẳng trên $S$, thì trong (5.1) khả
năng nâng địa phương được bảo đảm; mặt khác,
$\mathcal J/\mathcal J^2$ tự do địa phương trên $Y_0$ và ta có thể viết
\[
\mathcal A=\mathcal N_{Y_0/X_0}\otimes_{\mathcal O_{S_0}}\mathcal J,
\]
trong đó thừa số thứ nhất ở vế phải là bó pháp tuyến của $Y_0$ trong $X_0$.
Áp dụng tiêu chuẩn cơ bản về tính trơn ([3], III, 3.1), chẳng hạn ta thu được:

\result{HỆ QUẢ}{5.4} Trong các điều kiện của (5.3), giả sử rằng
$\mathcal F=\mathcal O_X$, rằng $X$ phẳng trên $S$, và rằng tiền lược đồ con
đóng $Y_0$ của $X_0$ tương ứng với $\mathcal G_0$ là một giao đầy đủ địa
phương. Khi đó, không gian tiếp tuyến Zariski của $Q_s$ tại $x$ đẳng cấu chính
tắc với $H^0(Y_0,\mathcal N_{Y_0/X_0})$. Nếu
$H^1(Y_0,\mathcal N_{Y_0/X_0})=0$, thì tiền lược đồ Hilbert
$\operatorname{Hilb}_{X/S}$ trơn trên $S$ tại điểm $x$.
($\mathcal N_{Y_0/X_0}$ là bó pháp tuyến của $Y_0$ trong $X_0$.)

\result{CHÚ GIẢI}{5.5} Phát biểu này áp dụng đặc biệt khi $Y_0$ là một giao
đầy đủ trong $X_0$ được xác định bởi một phương trình, tức là một ``ước
Cartier'' hiệu dụng. Khi đó, $\mathcal N_{Y_0/X_0}$ đẳng cấu với bó trên $Y_0$
cảm sinh bởi bó khả nghịch $\mathcal J^{-1}$ trên $X_0$ được xác định bởi
ước $Y_0$. Tình huống này xuất hiện đặc biệt khi nghiên cứu các họ ước hiệu
dụng trên một đa tạp xạ ảnh không kỳ dị $X_0$. Đẳng cấu giữa không gian tiếp
tuyến Zariski tại điểm $x$ của $Q$ (hoặc, nếu muốn, của tập mở $D$ của $Q$
tương ứng với các ước) và $H^0(Y_0,\mathcal N_{Y_0/X_0})$ từng được biết trong
hình học đại số cổ điển dưới tên ``đồng cấu đặc trưng'' từ không gian thứ nhất
đến không gian thứ hai. Đồng cấu này chỉ được xác định khi $x$ là một điểm trơn
của đa tạp tham số $T$ của một ``họ liên tục đầy đủ'' các ước; theo quan điểm
của chúng ta, đa tạp tham số ấy là một thành phần bất khả quy của lược đồ $D$
với cấu trúc rút gọn cảm sinh. Khi đó, không gian tiếp tuyến của $T$ tại $x$ là
một không gian con của không gian tiếp tuyến của $D$ tại $x$; do đó đồng cấu
đặc trưng cổ điển quả là đơn ánh, nhưng chỉ toàn ánh dưới các điều kiện bổ
sung, chẳng hạn nếu $D$ là nguyên tại $x$. Thật vậy, Zappa [8] đã xây dựng
một ví dụ
\src{272}
(với $X$ là một mặt xạ ảnh không kỳ dị trên trường số phức) trong đó, ngay cả
tại điểm tổng quát của $T$, đồng cấu đặc trưng cũng không toàn ánh. Do đó, điều
này có nghĩa là $D$ không nguyên ngay cả tại điểm tổng quát của thành phần bất
khả quy đang xét. Điều này cho thấy một cách đặc biệt nổi bật rằng các đa tạp
có phần tử lũy linh là cần thiết để hiểu những hiện tượng thuộc lý thuyết cổ
điển nhất về các mặt.


\subsection*{5.6}

Trong (5.4), ta đã nêu một tiêu chuẩn về tính trơn, áp dụng đặc biệt cho các
lược đồ ước. Trong [6], Kodaira đã nêu một tiêu chuẩn khác, cụ thể là sự triệt
tiêu của $H^1(X_0,\mathcal L)$, trong đó $\mathcal L=\mathcal J_0^{-1}$ là bó
khả nghịch trên $X_0$ được xác định bởi ước $Y_0$; tiêu chuẩn này đúng khi $S$
là phổ của một trường có đặc số 0, và được chứng minh trong [6] bằng phương
pháp siêu việt khi trường cơ sở là $\mathbf C$. Ta lưu ý rằng, nói chung, khi
$S$ lại là tùy ý, điều kiện của Kodaira là một điều kiện đủ để cấu xạ chính tắc
\[
D\longrightarrow\operatorname{Pic}_{X/S}
\]
từ tiền lược đồ các ước đến tiền lược đồ Picard của $X/S$ trơn tại điểm $x$
đang xét (điều này dễ dàng được kiểm tra bằng tiêu chuẩn thông thường về tính
trơn, một khi đã biết sự tồn tại của $\operatorname{Pic}_{X/S}$). Do đó, nếu
thêm vào đó $\operatorname{Pic}_{X/S}$ trơn trên $S$ tại điểm ảnh của $x$
(chẳng hạn nếu $\operatorname{Pic}_{X/S}$ trơn trên $S$), thì $D$ trơn trên
$S$ tại $x$. Mặt khác, Cartier đã chứng minh rằng mọi tiền lược đồ nhóm thuộc
kiểu hữu hạn địa phương trên một trường $k$ có đặc số 0 đều trơn trên $k$. Kết
hợp hai kết quả này, ta thu lại kết quả của Kodaira. Cũng cần lưu ý rằng từ các
nhận xét này suy ra rằng, trên một trường $K$ có đặc số $p>0$, khi
$\operatorname{Pic}_{X/S}$ không trơn trên $k$ (đó là trường hợp khi $X$ là
mặt Igusa), điều kiện $H^1(X_0,\mathcal L)=0$ trái lại kéo theo rằng $D$ không
trơn tại $x$, và thậm chí không rút gọn tại $x$ nếu $K$ đóng đại số.

Cuối cùng, ta còn nêu kết quả sau đây, vốn đóng một vai trò quan trọng trong
nghiên cứu vi phân các không gian phân thớ:

\result{MỆNH ĐỀ}{5.7} Cho $X$ là một tiền lược đồ hữu hạn và phẳng trên tiền
lược đồ Noether địa phương $S$, và cho $Z$ là một tiền lược đồ trên $S$ sao
cho $\Pi_{X/S}(Z/X)$ tồn tại (điều này đúng nếu $Z$ tựa xạ ảnh trên $X$). Khi
đó, nếu $Z$ trơn trên $X$, thì $\Pi_{X/S}(Z/X)$ trơn trên $S$.

Đây là một hệ quả trực tiếp của định nghĩa và của tiêu chuẩn thông thường về
tính trơn ([3], III, 3.1). Ta lưu ý rằng khi $X$ hữu hạn và phẳng trên $S$,
vấn đề về sự tồn tại của
$Z/K$ có thể
được xử lý rất sơ cấp, không cần
\src{273}
sử dụng lý thuyết lược đồ Hilbert. Chẳng hạn, ta thấy rằng nếu $X$ là
radiciel trên $S$, thì $\Pi_{X/S}(Z/X)$ tồn tại mà không có bất kỳ hạn chế nào
đối với $Z$. Một ví dụ khác, cho $T$ là một $S$-tiền lược đồ, và cho $T_n$ là
``lân cận vô cùng bé bậc $n$'' của đường chéo của $T\times_ST$ trong
$T\times_ST$, được trang bị các cấu xạ $p_1,p_2:T_n\longrightarrow T$ cảm
sinh bởi hai phép chiếu. Ta xem $T_n$ như một tiền lược đồ hữu hạn trên $T$
qua $p_1$, và giả sử thêm rằng nó phẳng trên $T$ (điều này đúng nếu $T$ trơn
trên $S$). Với mọi tiền lược đồ $X$ trên $T$, đặt
\[
(X/T/S)^{(n)}=
\Pi_{T_n/S}\bigl(p_2^*(X/T)/T_n\bigr),
\]

đó là một tiền lược đồ trên $T$ được gọi là \emph{phân thớ các mầm tiết diện
bậc $n$ của $X$ trên $T$} (tương đối với $S$). Nó phụ thuộc theo hàm tử vào
$X$, và trơn trên $T$ nếu $X$ cũng vậy.

\fsection{6}{Quan hệ với khái niệm chuẩn và các tích đối xứng}

Cho $S$ là một tiền lược đồ, $X$ và $Y$ là các $S$-tiền lược đồ,
\[
u:(X/S)^n\longrightarrow Y
\]
là một $S$-cấu xạ đối xứng từ lũy thừa Descartes thứ $n$ của $X/S$ vào $Y$.
Để đơn giản, ta giả sử $S$ là Noether địa phương, và $X$ cùng $Y$ thuộc kiểu
hữu hạn trên $S$. Khi đó, với mỗi môđun coherent $\mathcal F$ trên $X$, có giá
hữu hạn trên $S$, phẳng trên $S$ và có hạng tương đối trên $S$ bằng $n$ (tức
là $f_*(\mathcal F)$ là một môđun tự do địa phương hạng $n$ trên $S$), ta có
thể liên kết một cách tự nhiên một tiết diện của $Y$ trên $S$:
\[
\pi^u_{X/S}(\mathcal F)\in\Gamma(Y/S).
\]
Ở đây ta sẽ không đưa ra định nghĩa hình thức, mà chỉ lưu ý rằng hình thức
luận thu được là một sự tổng quát hóa tự nhiên của hình thức luận thông thường
về chuẩn và vết. Khi lũy thừa đối xứng thứ $n$ của $X$ trên $S$ tồn tại (chẳng
hạn khi các quỹ đạo của nhóm đối xứng $\mathfrak S_n$ tác động trên $(X/S)^n$
được chứa trong các tập mở affine), ta có thể lấy
$Y$ bằng $\operatorname{Symm}^n_S(X)$, và thu được một phần tử chính tắc
\[
\pi_{X/S}(\mathcal F)\in\Gamma(\operatorname{Symm}^n_S(X)/S),
\]
\src{274}
cho phép khôi phục các $\pi^u_{X/S}(\mathcal F)$. Một trường hợp quan trọng
khác là khi $X$ là một monoid giao hoán trên $S$, và $X=Y$, còn cấu xạ $u$
xuất phát từ luật hợp thành của $X$. Khi đó ta sẽ chỉ viết $\pi(\mathcal F)$
cho tiết diện của $X$ trên $S$ liên kết với môđun $\mathcal F$ trên $X$.

Bây giờ giả sử ta có một môđun coherent $\mathcal F$ trên $X$ sao cho
$\operatorname{Quot}_{\mathcal F/X/S}$ tồn tại, hoặc ít nhất sao cho hàm tử
$\underline{\operatorname{Quot}}^n_{\mathcal F/X/S}$ gán cho mỗi $S'$ trên
$S$ tập hợp các bó thương coherent $\mathcal M$ của
$\mathcal F'=\mathcal F\otimes_{\mathcal O_S}\mathcal O_{S'}$ phẳng trên
$S'$ và có hạng tương đối $n$, được biểu diễn bởi một $S$-tiền lược đồ
$\operatorname{Quot}^n_{\mathcal F/X/S}$. (Khi $X$ quasi-xạ ảnh trên $S$,
thì $\operatorname{Quot}^n_{\mathcal F/X/S}$ quả thực tồn tại và, với ký hiệu
của mục 3, chính là $\operatorname{Quot}^P_{\mathcal F/X/S}$, trong đó $P$ là
đa thức chỉ còn số hạng hằng $n$.) Vì phép dựng
$\pi^u_{X/S}(\mathcal M)$ tương thích với phép đổi cơ sở, ta thu được một cấu
xạ chính tắc
\[
\pi^u_{X/S}:\operatorname{Quot}^n_{\mathcal F/X/S}\longrightarrow Y,
\]
đặc biệt, nếu lũy thừa đối xứng thứ $n$ của $X$ trên $S$ tồn tại:
\[
\pi_{X/S}:\operatorname{Quot}^n_{\mathcal F/X/S}
\longrightarrow\operatorname{Symm}^n_S(X).
\]
Trường hợp quan trọng nhất là $\mathcal F=\mathcal O_X$, cho ta một cấu xạ:
\[
\pi_{X/S}:\operatorname{Hilb}^n_{X/S}
\longrightarrow\operatorname{Symm}^n_S(X).
\]
Hiển nhiên đây là một đẳng cấu khi $n=0$ hoặc $n=1$. Nhưng khi
$n\geq1$, ngay cả nếu $S$ là phổ của một trường $k$ và $X$ trơn trên
$S$, nói chung đây không phải là một đẳng cấu, thậm chí cũng không phải là một
cấu xạ đơn ánh, vì một lược đồ con chiều 0 của $X$ (chẳng hạn tương ứng với
một iđêan $\mathfrak i$ nguyên sơ đối với iđêan cực đại trong vành địa phương
$\mathcal O_{X,x}$, với $x$ là một điểm đóng của $X$) không được xác định khi
chỉ biết chu trình do nó xác định (trong trường hợp này, khi biết đối chiều
trên $k$ của $\mathfrak i$ trong $\mathcal O_{X,x}$). Ta chỉ có thể nói điều
sau đây ($S$ lại là tùy ý):

\medskip
\noindent a. Nếu $X$ trơn trên $S$, thì cấu xạ chuẩn xác định một đẳng cấu từ
tập mở của $\operatorname{Hilb}^n_{X/S}$ phân loại các phủ étale hạng $n$ nằm
trong $X$ (xem số \textsuperscript{o} 4, (b)) lên tập mở của
$\operatorname{Symm}^n_S(X)$ tương ứng với các $n$-chu trình không có thành
phần bội.
\src{275}

\medskip
\noindent b. Nếu thêm vào đó $X$ có chiều tương đối 1 trên $S$, thì cấu xạ
chuẩn thậm chí xác định một đẳng cấu từ $\operatorname{Hilb}^n_{X/S}$ lên
$\operatorname{Symm}^n_{X/S}$.

Sở dĩ như vậy là vì một lược đồ con chiều 0 của một đường cong đại số không
kỳ dị được xác định hoàn toàn bởi chu trình tương ứng. Cùng nhận xét ấy áp
dụng tổng quát hơn cho các ước Cartier hiệu dụng trên một lược đồ đại số không
kỳ dị (và không loại trừ khả năng rằng, trong trường hợp rất đặc biệt này, đa
tạp Chow cho cùng một đối tượng như lược đồ Hilbert).

\fsection{7}{Các bổ sung và câu hỏi}

Như J.-P. Serre đã nhận xét, từ một ví dụ nổi tiếng của Nagata suy ra rằng có
thể tìm được một lược đồ $S$, là phổ của một trường $k$, một $S$-lược đồ $S'$,
là phổ của một mở rộng bậc hai $k'$ của $k$, và cuối cùng một $S'$-lược đồ $X$
riêng và trơn (nhưng không xạ ảnh) chiều 3, sao cho $\Pi_{S'/S}(X/S')$ không
tồn tại. Điều này càng kéo theo rằng lược đồ Hilbert
$\operatorname{Hilb}^2_{X/S}$ không tồn tại (thậm chí cả $k$-lược đồ biểu diễn
các phủ étale hạng 2 của $S$ nằm trong $X$ cũng không tồn tại, và càng không có
bình phương đối xứng của $X$; xem mục trước). Do đó, khả năng thực hiện các
phép dựng không xạ ảnh trong hình học đại số bị đặt dưới những hạn chế nghiêm
trọng. (Tuy nhiên, có thể tin rằng những hạn chế như vậy sẽ không xuất hiện
trong hình học giải tích, cũng như chúng không xuất hiện trong hình học hình
thức (xem [2], II).) Mặt khác, nếu $X$ là một lược đồ riêng trên phổ $S$ của
một trường $k$, và nếu $Z$ tựa xạ ảnh trên $X$, thì $\Pi_{X/S}(Z/X)$ tồn tại
và là một lược đồ, hợp rời nhau của một dãy các lược đồ tựa xạ ảnh trên $S$
(như trong trường hợp xạ ảnh (3.1)). Thật vậy, ta quy về trường hợp bản thân
$X$ là xạ ảnh bằng cách để $X$ được trội bởi một $S$-lược đồ xạ ảnh $X'$; ở đây
chúng tôi không trình bày chi tiết chứng minh, vốn còn sử dụng kết quả về việc
phân tích một cấu xạ hữu hạn thành hợp được nêu trong [2], I, A,
n\textsuperscript{o} 2 (b). Sự thành công của phương pháp nằm ở chỗ, vì $S$ là
phổ của một trường, $X'$ xuất hiện trong bổ đề Chow tự động phẳng trên $S$.

Tôi không biết kết quả có còn đúng nếu không đặt giả thiết nào lên $S$ hay
không, khi chỉ giả sử $X$ riêng và phẳng trên $S$, còn $Z$ tựa xạ ảnh trên
$X$. Một trường hợp quan trọng trong ứng dụng là khi $Z$ là một lược đồ con
đóng của $X$; nếu khi ấy $\Pi_{X/S}(Z/X)$ tồn tại, thì nó nhất thiết là một
lược đồ con đóng của $S$. Ta có thể
\src{276}
dựng nó trực tiếp khá đơn giản khi $X$ xạ ảnh trên $S$, mà không dùng lý
thuyết lược đồ Hilbert; tổng quát hơn, phương pháp được sử dụng cho thấy rằng
nếu $Z$ affine trên $X$, thì $\Pi_{X/S}(Z/X)$ tồn tại và affine trên $S$. Nó
cũng cho thấy rằng nếu $X$ riêng và phẳng trên $S$ (không nhất thiết xạ ảnh
trên $S$), thì với mọi phân thớ vectơ $Z$ tầm thường địa phương trên $X$,
$\Pi_{X/S}(Z/X)$ tồn tại và là một phân thớ vectơ trên $S$. Rất nên xem xét
lại và thống nhất các kết quả này.

\section*{TÀI LIỆU THAM KHẢO}
\addcontentsline{toc}{section}{Tài liệu tham khảo}

\begin{description}
\item[{[1]}] BOREL (A.) và SERRE (J.-P.). --- Le théorème de Riemann--Roch,
\emph{Bull. Soc. math. France}, tập 86, 1958, tr. 97--136.

\item[{[2]}] GROTHENDIECK (Alexander). --- Technique de construction et
théorèmes d'existence en géométrie algébrique, I, II, III, \emph{Séminaire
Bourbaki}, tập 13, 1959/60, n\textsuperscript{o} 190, 29 tr.,
n\textsuperscript{o} 195, 22 tr.; và tập 14, 1960/61,
n\textsuperscript{o} 212, 20 tr.

\item[{[3]}] GROTHENDIECK (Alexander). --- Séminaire de Géométrie algébrique,
I, II, III, IV. --- Paris, Institut des hautes Études scientifiques, 1960/61
(in rô-nê-ô).

\item[{[4]}] GROTHENDIECK (Alexander). --- Technique de construction en
géométrie analytique, IV, V, \emph{Séminaire Cartan}, tập 13, 1960/61,
n\textsuperscript{o} 11 và 12.

\item[{[5]}] GROTHENDIECK (A.) và DIEUDONNÉ (J.). --- \emph{Éléments de
géométrie algébrique}, I: Le langage des schémas. --- Paris, Presses
universitaires de France, 1960 (Institut des hautes Études scientifiques,
Publications mathématiques, 4); II, III, IV (sắp xuất bản).

\item[{[6]}] KODAIRA (K.). --- Characteristic linear systems of complete
continuous systems, \emph{Amer. J. of Math.}, tập 78, 1956, tr. 716--744.

\item[{[7]}] SERRE (Jean-Pierre). --- Exemples de variétés projectives en
caractéristique $p$ non relevables en caractéristique 0, \emph{Proc. Nat. Acad.
Sc. U.S.A.}, tập 47, 1961, tr. 108--109.

\item[{[8]}] ZAPPA (G.). --- Sull'esistenza, sopra le superficie algebriche, di
sistemi continui completi infiniti, la cui curva è a serie caratteristica
incompleta, \emph{Pont. Acad. Sc. Acta}, tập 9, 1945, tr. 91--93.
\end{description}

\section*{BỔ SUNG}
\addcontentsline{toc}{section}{Bổ sung}
\begin{center}
{\small [bổ sung khi sửa bản in thử]}
\end{center}

Giờ đây có thể thấy rằng các phỏng đoán của [2], III,
n\textsuperscript{o} 8 là sai, ngay cả đối với các đa tạp không kỳ dị trên một
trường đặc trưng 0, cả về sự tồn tại lẫn tính tựa xạ ảnh của thương, và ngay cả
khi $G$ tác động với đồ thị đóng.
